% =====================================================
\textbf{UC-KON-01: Opracowanie i zatwierdzenie specyfikacji pojazdu}

\begin{longtable}{|p{4cm}|p{11cm}|}
\hline
\textbf{Element} & \textbf{Opis} \\
\hline

Identyfikator & UC-KON-01 \\
\hline

Nazwa & Opracowanie i zatwierdzenie specyfikacji pojazdu \\
\hline

Opis & 
Proces wyboru parametrów pojazdu oparty na predefiniowanych pakietach wyposażenia, silnikach, skrzyniach biegów oraz kolorach nadwozia, zakończony zatwierdzeniem konfiguracji. \\
\hline

Aktor główny & Klient \\
\hline

Aktorzy pomocniczy & - \\
\hline

Warunki wstępne & 
Kontekst Katalog i Konfigurator otrzymał zdarzenie \textit{InitiateConfiguratorSession}. \\
\hline

Warunki końcowe & 
Zatwierdzona specyfikacja zostaje zapisana w lokalnej bazie danych konfiguratora, a system emituje zdarzenie \textit{SpecificationCompleted}. \\
\hline

Scenariusz główny & 
\begin{enumerate}
    \item System otwiera sesję konfiguratora i prezentuje Klientowi interfejs z modelami samochodów, wersjami wyposażenia (np. pakiet A1, A2), silnikami (np. B1, B2), rodzajami skrzyni biegów (np. C1, C2) oraz kolorami.
    \item Klient wybiera jeden główny pakiet wyposażenia (np. A1).
    \item Klient dobiera do pakietu silnik (np. B1), rodzaj skrzyni biegów (np. C1) oraz kolor.
    \item System na bieżąco weryfikuje proste reguły wykluczające (np. silnik B2 nie może być wybrany ze skrzynią biegów C1).
    \item Klient zatwierdza specyfikację.
    \item System weryfikuje ostateczną kompletność konfiguracji i kalkuluje cenę.
    \item System zapisuje specyfikację i emituje zdarzenie \textit{SpecificationCompleted}.
\end{enumerate} \\
\hline

Scenariusze alternatywne & 
\begin{itemize}
    \item [A1] Niedozwolona kombinacja: W kroku 4 system wykrywa, że wybrana kombinacja jest zablokowana przez producenta (np. silnik B2 nie może być wybrany ze skrzynią C1). System blokuje możliwość zatwierdzenia i prosi Klienta o zmianę silnika, skrzyni biegów lub pakietu. Kontekst nie emituje zdarzenia końcowego.
    \item [A2] Przerwanie sesji: Klient opuszcza konfigurator przed zatwierdzeniem. System zapisuje konfigurację jako wersję roboczą.
\end{itemize} \\
\hline
\end{longtable}

% =====================================================
\textbf{UC-KON-02: Automatyczna aktualizacja cennika i katalogu}

\begin{longtable}{|p{4cm}|p{11cm}|}
\hline
\textbf{Element} & \textbf{Opis} \\
\hline

Identyfikator & UC-KON-02 \\
\hline

Nazwa & Automatyczna aktualizacja cennika i katalogu \\
\hline

Opis & 
Zautomatyzowany proces pobierania, walidacji i zapisu nowego katalogu modeli oraz cennika, inicjowany przez zewnętrzny system producenta/importera (Blackbox). \\
\hline

Aktor główny & Zewnętrzny system producenta/importera \\
\hline

Aktorzy pomocniczy & - \\
\hline

Warunki wstępne & 
System zewnętrzny importera/producenta udostępnia nowy pakiet danych katalogowych (np. poprzez webhook, szynę zdarzeń lub API). \\
\hline

Warunki końcowe & 
Lokalna baza konfiguratora zostaje zaktualizowana, a kontekst emituje zdarzenie \textit{CatalogUpdated}. \\
\hline

Scenariusz główny & 
\begin{enumerate}
    \item Kontekst odbiera sygnał/dane o nowym cenniku z systemu Blackbox (np. poprzez zdarzenie integracyjne lub żądanie API).
    \item System parsuje otrzymany pakiet danych, wykorzystując warstwę translacji (ACL), aby dostosować format zewnętrzny do modelu lokalnego.
    \item System przeprowadza walidację strukturalną oraz logiczną (weryfikacja spójności reguł wyposażenia i poprawności cen).
    \item System zapisuje nowy katalog w lokalnej bazie danych, zastępując bieżące dane (poprzednie zostają oznaczone jako archiwalne).
    \item Kontekst emituje na szynę danych zdarzenie \textit{CatalogUpdated}.
\end{enumerate} \\
\hline

Scenariusze alternatywne & 
\begin{itemize}
    \item [A1] Błąd walidacji lub translacji danych: W kroku 2 lub 3 system wykrywa błędy krytyczne (np. brak cen, niezgodny format). System przerywa aktualizację, odrzuca pakiet i emituje techniczne zdarzenie o błędzie integracji (np. \textit{CatalogUpdateFailed}), logując szczegóły dla wsparcia IT.
\end{itemize} \\
\hline
\end{longtable}
